\subsection{Évaluation par Couverture dans un Arbre à Deux Périodes}

Pour calculer la valeur de l'option aujourd'hui, nous devons procéder à l'envers (calcul à rebours). Puisque le marché est complet à chaque nœud, nous pouvons appliquer notre argument de couverture locale par induction rétrograde. Nous calculons d'abord la valeur de l'option à la mi-journée (nœuds D et E) en utilisant les flux finaux (nœuds A, B et C). Une fois ces valeurs intermédiaires connues, nous recommençons l'opération pour trouver la valeur initiale à t=0.

Considérons notre option d'achat (Call) de prix d'exercice K = 100. Son profil de flux à l'échéance (t=1) est : 10 en A, 0 en C, 0 en B. Le taux d'intérêt sans risque est toujours supposé nul.

\begin{enumerate}[label=Étape \arabic* :, leftmargin=*]
    
    \item Évaluation au nœud E (L'action vaut 95 à t=0,5)
    \begin{itemize}[label=$\cdot$]
        \item Si nous sommes au nœud E, l'action finira à l'échéance soit à 100 (État C), soit à 90 (État B). 
        \item Dans les deux cas, le Call de Strike 100 ne vaudra strictement rien (0).
        \item Conclusion immédiate : La valeur de l'option au nœud E est Opt(E) = 0.
    \end{itemize}
    
    \vspace{0.2cm}
    \item Évaluation au nœud D (L'action vaut 105 à t=0,5)
    \begin{itemize}[label=$\cdot$]
        \item À partir de D, l'action peut monter à 110 (le Call vaudra 10) ou baisser à 100 (le Call vaudra 0).
        \item Construisons le portefeuille de couverture : nous achetons une quantité delta d'actions.
        \item Valeur future du portefeuille : 110 delta - 10 (si hausse) et 100 delta - 0 (si baisse).
        \item Pour éliminer le risque, nous égalisons : 110 delta - 10 = 100 delta, ce qui donne 10 delta = 10, soit delta = 1.
        \item Avec 1 action, la valeur garantie à t=1 est de 100.
        \item Par arbitrage, le portefeuille au nœud D vaut 100 : 105(1) - Opt(D) = 100, ce qui implique Opt(D) = 5.
    \end{itemize}

    \vspace{0.2cm}
    \item Évaluation au nœud initial (L'action vaut 100 à t=0)
    \begin{itemize}[label=$\cdot$]
        \item Nous connaissons désormais les valeurs à la mi-journée : 5 si l'action monte à 105, 0 si elle baisse à 95.
        \item Construisons le portefeuille initial : nous achetons une nouvelle quantité delta d'actions.
        \item Valeur future (à t=0,5) : 105 delta - 5 (si hausse vers D) et 95 delta - 0 (si baisse vers E).
        \item Égalisation des flux : 105 delta - 5 = 95 delta, ce qui donne 10 delta = 5, soit delta = 0,5.
        \item Avec 0,5 action, la valeur garantie à t=0,5 est de 95(0,5) = 47,5.
        \item Par arbitrage, la valeur initiale du portefeuille aujourd'hui doit être de 47,5 : 100(0,5) - Opt(0) = 47,5.
        \item Le prix final de l'option aujourd'hui est Opt(0) = 2,5.
    \end{itemize}

\end{enumerate}

Cet exemple démontre mathématiquement la notion de couverture dynamique. Remarquez que le delta (la quantité d'actions à détenir pour se couvrir) n'est pas constant. Il vaut 0,5 au début de la journée. Si l'action monte à 105, le vendeur de l'option devra ajuster son portefeuille en achetant 0,5 action supplémentaire pour que son delta passe à 1. La couverture parfaite exige un ajustement à chaque pas de temps.

\subsection{Évaluation Risque-Neutre dans un Arbre à Deux Périodes}

Bien que l'argument de couverture nous garantisse un prix mathématiquement parfait, les calculs de portefeuilles deviennent vite fastidieux. L'approche par la probabilité risque-neutre, bien que plus abstraite, est beaucoup plus simple à appliquer. Elle donnera toujours un résultat strictement identique à la méthode par couverture. Il existe deux manières de l'utiliser : en remontant l'arbre pas à pas (induction rétrograde), ou en calculant directement les probabilités globales d'atteindre chaque état final.

Méthode 1 : Calcul pas à pas (Induction rétrograde)

Nous reprenons le même arbre. Le taux d'intérêt est nul, donc l'espérance du prix futur de l'action doit égaler son prix actuel à chaque nœud.

\begin{enumerate}[label=Étape \arabic* :, leftmargin=*]
    \item Évaluation au nœud D (L'action vaut 105)
    \begin{itemize}[label=$\cdot$]
        \item L'action peut passer à 110 ou 100. La probabilité p (Hausse) doit vérifier : 110p + 100(1-p) = 105.
        \item La seule probabilité valide est p = 1/2.
        \item La valeur de l'option est son espérance : Opt(D) = 0,5(10) + 0,5(0) = 5.
    \end{itemize}

    \vspace{0.2cm}
    \item Évaluation au nœud E (L'action vaut 95)
    \begin{itemize}[label=$\cdot$]
        \item Les flux finaux de l'option sont nuls dans tous les scénarios (États C et B).
        \item L'espérance est de zéro quelle que soit la probabilité. Opt(E) = 0.
    \end{itemize}

    \vspace{0.2cm}
    \item Évaluation initiale (L'action vaut 100)
    \begin{itemize}[label=$\cdot$]
        \item L'action passe à 105 ou 95. Pour que l'espérance soit 100, la probabilité de Hausse doit de nouveau être p = 1/2.
        \item La valeur initiale de l'option est l'espérance de ses valeurs futures intermédiaires (Opt(D) et Opt(E)).
        \item Opt(0) = 0,5(5) + 0,5(0) = 2,5.
    \end{itemize}
\end{enumerate}
Le résultat est identique au prix d'arbitrage obtenu par la méthode de couverture en delta.

Méthode 2 : Calcul par les probabilités globales des chemins

Plutôt que de calculer les valeurs intermédiaires de l'option, nous pouvons calculer directement la probabilité risque-neutre de nous retrouver dans chaque état final à t=1, puis calculer l'espérance finale en une seule fois. Sachant qu'à chaque nœud la probabilité de Hausse est de 1/2 :

\begin{itemize}[label=$\cdot$]
    \item État A (Action à 110) : Il faut 2 hausses successives. Probabilité = 1/2 fois 1/2 = 1/4.
    \item État C (Action à 100) : Il faut 1 hausse et 1 baisse (ou inversement). Probabilité = 1/4 + 1/4 = 1/2.
    \item État B (Action à 90) : Il faut 2 baisses successives. Probabilité = 1/2 fois 1/2 = 1/4.
\end{itemize}

Calcul direct de l'espérance de l'option aujourd'hui :
Opt(0) = 1/4(10) + 1/2(0) + 1/4(0) = 10/4 = 2,5.

Cette méthode est extrêmement puissante car elle permet de sauter toutes les étapes intermédiaires de calcul.

L'unique caractéristique qualitative cruciale de cet argument est qu'à chaque pas de temps, l'actif ne peut évoluer que vers deux nouvelles valeurs possibles. C'est ce qui garantit la complétude locale du marché. Le nombre total de pas de temps n'a aucune importance mathématique. On peut subdiviser le temps en autant d'étapes que l'on souhaite (10, 100, 1000 étapes). Tant que l'arbre reste binomial à chaque nœud, le prix d'arbitrage sera unique. L'interprétation par la réplication reste valide : le prix de l'option est exactement la somme d'argent qu'il faut investir aujourd'hui pour fabriquer un clone parfait. Ce clone nécessitant un rééquilibrage continu, la quantité d'actions à détenir (le Delta) et le montant de cash changent à chaque nœud de l'arbre en fonction du temps et du prix de l'action.

\subsection{Généralisation à N Périodes}

Il n'y a rien de magique dans un modèle à deux étapes. Tant que chaque nœud de l'arbre se divise exactement en deux branches, la logique d'induction rétrograde garantit un prix sans arbitrage unique. Au lieu de calculer à rebours (nœud par nœud, couche par couche) depuis l'échéance finale jusqu'à aujourd'hui, nous pouvons directement calculer les probabilités globales d'atteindre chaque nœud final de l'arbre, puis faire l'espérance mathématique des flux finaux.

Exemple d'un arbre à 4 pas de temps (Taux d'intérêt = 0) :
Imaginons un arbre à 4 étapes, où le cours de l'action monte ou baisse de 2,5 à chaque pas. L'action part de 100 à t=0. 
Puisque le taux d'intérêt est nul, la probabilité risque-neutre de hausse (p) est obligatoirement de 0,5 à chaque nœud pour maintenir l'espérance locale.
Calculons la probabilité risque-neutre globale de chaque état final à la 4ème étape :

\begin{itemize}[label=$\cdot$]
    \item Arrivée à 110 : Il faut exactement 4 hausses successives.
    Il n'y a qu'un seul chemin. Probabilité : $1 \times (1/2)^4 = 1/16$.
    
    \item Arrivée à 105 : Il faut exactement 3 hausses et 1 baisse.
    Le nombre de combinaisons possibles est donné par le coefficient binomial $C_4^1 = 4$.
    Probabilité : $4 \times (1/2)^4 = 4/16 = 1/4$.
    
    \item Arrivée à 100 : Il faut exactement 2 hausses et 2 baisses.
    Nombre de combinaisons $C_4^2 = 6$.
    Probabilité : $6 \times (1/2)^4 = 6/16 = 3/8$.
    
    \item Arrivée à 95 : Il faut exactement 1 hausse et 3 baisses.
    Nombre de combinaisons $C_4^3 = 4$.
    Probabilité : $4 \times (1/2)^4 = 4/16 = 1/4$.
    
    \item Arrivée à 90 : Il faut exactement 0 hausse et 4 baisses.
    Il n'y a qu'un seul chemin. Probabilité : $1 \times (1/2)^4 = 1/16$.
\end{itemize}

Tarification d'un Call de Strike 100 :
Nous multiplions la probabilité de chaque état par la valeur intrinsèque du Call dans cet état.
Opt(K=100) = (1/16) fois (110 - 100) + (1/4) fois (105 - 100) + 0 + 0 + 0
Opt(K=100) = 10/16 + 5/4 = 10/16 + 20/16 = 30/16 = 15/8.

\vspace{0.5cm}
\begin{center}
\begin{tikzpicture}[>=stealth, sloped]
    % Style pour les mini noeuds d'un grand arbre
    \tikzstyle{mini_node} = [circle, draw=vscodeBlue, fill=codeBackground, thick, minimum size=0.4cm, inner sep=0pt, font=\tiny]
    \tikzstyle{time_style} = [font=\small\color{gray}]

    % Dimensions de l'arbre
    \def\N{5} % Nombre d'étapes
    \def\dx{1.8} % Espace horizontal entre les étapes
    \def\dy{0.8} % Espace vertical entre les noeuds

    % Lignes de temps
    \foreach \i in {0,...,\N} {
        \node[time_style] at (\i*\dx, \N*\dy*0.5 + 0.8) {$t=\i$};
    }

    % Génération automatique des noeuds et des liens (Arbre Recombinant)
    \foreach \i in {0,...,\N} {
        \foreach \j in {0,...,\i} {
            % Calcul de la position Y : Centré autour de Y=0
            \pgfmathsetmacro{\ypos}{(\i - 2*\j)*\dy*0.5}
            \node[mini_node] (N-\i-\j) at (\i*\dx, \ypos) {};
            
            % Dessiner les liens vers la génération précédente (si on n'est pas à i=0)
            \ifnum\i>0
                \pgfmathtruncatemacro{\prevI}{\i-1} % CORRECTION ICI
                % Lien depuis le noeud supérieur (baisse)
                \ifnum\j>0
                    \pgfmathtruncatemacro{\prevJup}{\j-1} % CORRECTION ICI
                    \draw[->, vscodeBlue!60] (N-\prevI-\prevJup) -- (N-\i-\j);
                \fi
                % Lien depuis le noeud inférieur (hausse)
                \ifnum\j<\i
                    \pgfmathtruncatemacro{\prevJdown}{\j} % CORRECTION ICI
                    \draw[->, vscodeBlue!60] (N-\prevI-\prevJdown) -- (N-\i-\j);
                \fi
            \fi
        }
    }
    
    % Étiquette pour montrer la multiplication des chemins
    \node[right=0.2cm of N-\N-0, font=\scriptsize] {5 Hausses};
    \node[right=0.2cm of N-\N-5, font=\scriptsize] {5 Baisses};
    \node[right=0.2cm of N-\N-2, font=\scriptsize] {Chemins multiples};
    
\end{tikzpicture}
\end{center}
\vspace{0.5cm}

\begin{theorembox}[La Formule Globale pour N pas de temps]
Considérons une action de prix initial S. L'arbre compte $N$ étapes. À chaque étape, l'action monte ou baisse d'un montant absolu $x$.
À la fin de l'arbre, la valeur de l'action dépend uniquement du nombre total de hausses (noté $j$). Si l'action a fait $j$ hausses, elle a obligatoirement fait $(N - j)$ baisses.
Le prix final de l'action $S_j$ dans l'état avec $j$ hausses s'écrit :
\begin{equation}
    \label{eq:fsec3-prix_sj}
    S_j = S + j(x) - (N - j)(x) = S + (2j - N)x
\end{equation}

Puisque la probabilité d'une hausse est de $1/2$ (avec un taux nul), la probabilité totale d'atteindre ce prix final $S_j$ suit une distribution binomiale :
\begin{equation}
    \label{eq:fsec3-prob_binom}
    Prob(S_j) = \frac{1}{2^N} \binom{N}{j}
\end{equation}
Où $\binom{N}{j}$ représente le nombre de chemins différents permettant d'obtenir $j$ hausses sur $N$ lancers.

Par conséquent, la valeur actuelle d'un produit dérivé qui génère un flux $f(S)$ à l'échéance est simplement l'espérance de tous les flux terminaux pondérés par cette probabilité :
\begin{equation}
    \label{eq:fsec3-prix_n_pas}
    Opt_0 = \frac{1}{2^N} \sum_{j=0}^{N} \binom{N}{j} f\left(S + (2j - N)x\right)
\end{equation}
\end{theorembox}

L'objectif ultime de cette modélisation mathématique n'est pas de calculer un arbre géant à la main, mais de comprendre ce qui se passe mathématiquement lorsque l'on fait tendre le nombre d'étapes $N$ vers l'infini. Pour que la limite ait un sens mathématique et économique, il faudra réduire la taille de la variation de prix ($x$) au fur et à mesure que le nombre d'étapes ($N$) augmente, tout en préservant certaines propriétés de la volatilité de l'action. C'est précisément ce passage à la limite qui nous fera basculer de l'arbre binomial discret au mouvement brownien continu, fondement du modèle de Black-Scholes.

\subsection{Le Modèle Normal}

Pour obtenir un prix d'option infiniment précis, il faut découper le temps en intervalles de plus en plus petits, c'est-à-dire faire tendre le nombre de pas de temps vers l'infini. Cependant, nous ne pouvons pas conserver une variation de prix constante à chaque étape, sous peine de voir la volatilité finale exploser. Il faut que la taille du saut diminue à mesure que le nombre d'étapes augmente. La méthode exacte pour réduire cette taille définit la nature du modèle continu que nous obtiendrons à la limite.

Considérons un modèle simplifié où le taux d'intérêt est nul. L'actif a une valeur initiale $S_0$. L'échéance de l'option est $T$.
Nous imposons deux contraintes statistiques fondamentales sur le prix de l'actif à l'échéance $T$ :
\begin{itemize}[label=$\rightarrow$]
    \item L'espérance du prix final doit être égale au prix initial : $\mathbb{E}[S_T] = S_0$.
    \item La variance totale du prix final doit croître linéairement avec le temps : $Var(S_T) = \sigma^2 T$, où $\sigma$ représente la volatilité absolue.
\end{itemize}

Nous divisons la durée totale $T$ en $N$ pas de temps réguliers. Puisque les mouvements à chaque pas sont indépendants, la variance totale est la somme des variances locales. Pour obtenir une variance totale de $\sigma^2 T$ après $N$ pas, la variance de chaque micro-étape doit obligatoirement être de $\sigma^2 T / N$.

Supposons qu'à chaque étape, l'actif monte d'un montant $+x_N$ ou baisse d'un montant $-x_N$ avec une probabilité équiprobable de $0,5$.
L'espérance locale de ce mouvement est bien nulle : $0,5(x_N) + 0,5(-x_N) = 0$.
La variance locale de ce mouvement est l'espérance des carrés : $0,5(x_N)^2 + 0,5(-x_N)^2 = (x_N)^2$.

En égalisant cette variance locale avec notre cible, nous déduisons la taille exacte du saut par pas de temps :
\begin{equation}
    \label{eq:fsec3-taille_saut}
    (x_N)^2 = \frac{\sigma^2 T}{N} \implies x_N = \sigma \sqrt{\frac{T}{N}}
\end{equation}
Le mouvement de l'actif à chaque étape doit donc être proportionnel à la racine carrée du pas de temps.

Pour formaliser cette dynamique, introduisons une suite de variables aléatoires indépendantes $Z_i$. À l'étape $i$, $Z_i$ prend la valeur $+1$ ou $-1$, chacune avec une probabilité de $0,5$. Après $N$ étapes de temps, le prix de l'actif est la somme de son prix initial et de toutes ses variations :
\begin{equation}
    \label{eq:fsec3-prix_sn_discret}
    S_N = S_0 + \sigma \sqrt{\frac{T}{N}} \sum_{i=1}^{N} Z_i
\end{equation}
Le prix sans arbitrage de l'option d'achat (qui verse un flux final $f(S_T)$) est l'espérance mathématique calculée sur l'ensemble de ces chemins :
\begin{equation}
    \label{eq:fsec3-opt_sn_discret}
    Opt_0 = \mathbb{E}\left[ f\left( S_0 + \sigma \sqrt{\frac{T}{N}} \sum_{i=1}^{N} Z_i \right) \right]
\end{equation}

Le but de cette construction est d'étudier la limite de cette espérance lorsque $N$ tend vers l'infini. Dans l'équation ci-dessus, le terme $\sum_{i=1}^{N} Z_i$ représente une somme de variables aléatoires indépendantes et identiquement distribuées, de moyenne nulle. Le Théorème Central Limite affirme que lorsqu'on normalise une telle somme par la racine carrée de $N$, sa distribution de probabilité converge irrémédiablement vers une loi Normale standard (la fameuse courbe en cloche de Gauss). Par conséquent, à la limite continue, le prix futur de l'action $S_T$ suivra une distribution Normale. Ce modèle arithmétique constitue les bases historiques posées par Louis Bachelier en 1900.

\subsection{Convergence vers l'Intégrale de Valorisation}

\begin{theorembox}[Théorème Central Limite (TCL)]
Soit $X_1, X_2, \dots$ une suite de variables aléatoires indépendantes et identiquement distribuées (i.i.d.), ayant une moyenne finie $\mu$ et une variance finie non nulle $\sigma^2$.
Soit $S_n$ la somme de ces $n$ variables :
\begin{equation}
    \label{eq:fsec3-tcl_somme}
    S_n = X_1 + X_2 + \dots + X_n
\end{equation}
Alors, la distribution de la variable centrée et réduite :
\begin{equation}
    \label{eq:fsec3-tcl_normale}
    \frac{S_n - n\mu}{\sigma \sqrt{n}}
\end{equation}
converge vers celle d'une variable aléatoire Gaussienne de moyenne 0 et de variance 1 lorsque $n$ tend vers l'infini.
\end{theorembox}

Dans notre arbre binomial calibré, les variables aléatoires $Z_i$ prennent les valeurs $+1$ ou $-1$ avec probabilité $0,5$. Elles sont indépendantes, de moyenne nulle ($\mu = 0$) et de variance unitaire ($\sigma^2 = 1$). En appliquant le Théorème Central Limite, la somme de ces variables divisée par la racine du nombre de pas $N$ :
\begin{equation}
    \label{eq:fsec3-tcl_app}
    \frac{1}{\sqrt{N}} \sum_{i=1}^{N} Z_i
\end{equation}
converge vers une distribution Gaussienne standard, notée $\mathcal{N}(0,1)$, de moyenne 0 et de variance 1, lorsque $N \to \infty$.

Le prix de l'actif après $N$ étapes était donné par :
\begin{equation}
    \label{eq:fsec3-sn_limite}
    S_N = S_0 + \sigma \sqrt{T} \left( \frac{1}{\sqrt{N}} \sum_{i=1}^{N} Z_i \right)
\end{equation}
En passant à la limite continue ($N \to \infty$), le terme entre parenthèses devient la variable aléatoire $\mathcal{N}(0,1)$. La distribution du prix futur de l'actif converge donc vers :
\begin{equation}
    \label{eq:fsec3-st_continu}
    S_T = S_0 + \sigma \sqrt{T} \mathcal{N}(0,1)
\end{equation}

Le prix d'arbitrage de l'option est l'espérance mathématique de son flux final $f(S_T)$. Puisque $S_T$ suit une loi Normale, cette espérance se calcule formellement à l'aide de la fonction de densité de la loi Normale.
Le prix converge vers l'intégrale suivante :
\begin{equation}
    \label{eq:fsec3-integrale_continue}
    Opt_0 = \mathbb{E}[f(S_T)] = \frac{1}{\sqrt{2\pi}} \int_{-\infty}^{+\infty} f(S_0 + \sigma \sqrt{T} x) e^{-\frac{x^2}{2}} dx
\end{equation}
Pour certains types de flux finaux, comme celui d'une option d'achat classique vanille où $f(S) = \max(S - K, 0)$, cette intégrale peut être résolue analytiquement pour obtenir une formule de prix explicite.

La formule obtenue ici (le modèle Normal) suppose que le prix de l'actif suit une loi Normale (arithmétique). Bien qu'élégante mathématiquement, cette modélisation pose un problème financier majeur : une variable régie par une loi Normale peut prendre des valeurs négatives. Or, le prix d'une action, du fait de la responsabilité limitée des actionnaires, ne peut jamais tomber sous zéro. C'est cette aberration économique qui a conduit les chercheurs à abandonner le modèle Normal de Bachelier au profit du modèle Lognormal (géométrique) pour construire la théorie moderne de Black-Scholes.